\documentclass[conference]{IEEEtran}
\IEEEoverridecommandlockouts
\usepackage{cite}
\usepackage{amsmath,amssymb,amsfonts}
\usepackage{algorithmic}
\usepackage{graphicx}
\usepackage{textcomp}
\usepackage{xcolor}
\def\BibTeX{{\rm B\kern-.05em{\sc i\kern-.025em b}\kern-.08em
    T\kern-.1667em\lower.7ex\hbox{E}\kern-.125emX}}
\begin{document}

\title{Research Proposal: Enhancing AIOps with Reinforcement Learning through OpenR Framework Integration}

\author{\IEEEauthorblockN{1\textsuperscript{st} Wentao Ye}
    % \IEEEauthorblockA{\textit{dept. name of organization (of Aff.)} \\
    %     \textit{name of organization (of Aff.)}\\
    %     City, Country \\
    %     email address or ORCID}
}

\date{\today}

\maketitle

\begin{abstract}
    Managing complex cloud environments in IT Operations (ITOps) demands intelligent systems capable of emulating human reasoning to detect, diagnose, and resolve issues efficiently. 
    AIOpsLab structures IT operations into four key stages: Incident Detection, Triage after Localization, Root Cause Diagnosis, and Mitigation Resolution.
    Each stage requires iterative interaction with the cloud environment.

    The OpenR framework offers a novel reinforcement learning (RL) suite designed to model reasoning-intensive tasks by structuring reasoning processes as a sequence of steps (Q $\rightarrow$ \{R\} $\rightarrow$ A) within a Markov Decision Process (MDP). 
    
    This proposal explores integrating OpenR into AIOpsLab to develop an RL-driven solution for IT operations, leveraging OpenR's reasoning-focused approach to enhance performance across all operational stages.
\end{abstract}

\section{Background and Motivation}

\subsection{IT Operations Framework (AIOpsLab)}
AIOpsLab structures IT operations into four stages:
\begin{enumerate}
    \item \textbf{Incident Detection}: Identifying anomalies or issues in the system.
    \item \textbf{Triage after Localization}: Narrowing down potential causes and prioritizing them.
    \item \textbf{Root Cause Diagnosis}: Pinpointing the exact cause of the issue.
    \item \textbf{Mitigation Resolution}: Implementing corrective actions to resolve the issue.
\end{enumerate}

Each stage involves iterative reasoning and interaction with the cloud environment, resembling human problem-solving. An RL agent capable of emulating this reasoning can significantly enhance the speed, accuracy, and efficiency of IT operations.

\subsection{The OpenR Framework}
The OpenR framework proposes a structured methodology for reasoning tasks, particularly in scenarios that require iterative problem-solving and question-answering. It defines the reasoning process as a sequence (Q $\rightarrow$ \{R\} $\rightarrow$ A), where:
\begin{itemize}
    \item \textbf{Q}: The initiating question or prompt initiating the reasoning process.
    \item \textbf{R}: Intermediate reasoning steps taken to build toward a solution.
    \item \textbf{A}: Final answer or solution derived from the reasoning steps.
\end{itemize}

By framing reasoning as a Markov Decision Process (MDP), OpenR provides flexibility in modeling both sequential and branching reasoning paths:
\begin{itemize}
    \item \textbf{States} represent the sequence of reasoning steps taken so far.
    \item \textbf{Actions} involve selecting the next reasoning step or final answer.
    \item \textbf{Policy} refers to the LLM's function for generating reasoning steps and answers.
    \item \textbf{Process Reward Model (PRM)} provides feedback on the quality of reasoning steps and final answers, guiding the LLM policy to optimize its output.
\end{itemize}

This combination of sequential and branching reasoning allows OpenR to explore diverse solutions while maintaining logical consistency, making it a robust framework for reasoning-intensive tasks.

\subsection{Motivation for Integration}
Both OpenR and AIOpsLab emphasize the importance of reasoning in task-solving. 
Prior research in the math-solving domain has demonstrated that integrating reasoning frameworks like OpenR can significantly enhance AI agents' reasoning abilities over inference time. 
By performing a sequence of reasoning steps, agents exhibit improved problem-solving efficiency and accuracy. 
Inspired by these successes, integrating OpenR into AIOpsLab aims to:
\begin{enumerate}
    \item \textbf{Enhance Agent Performance}: Utilize OpenR's structured reasoning to improve AIOpsLab agents' capabilities in each operational stage.
    \item \textbf{Enable Adaptive Learning}: Implement OpenR's RL methodology to continuously refine AIOpsLab agents' decision-making during cloud interactions.
\end{enumerate}

\section{Objectives}
\begin{enumerate}
    \item \textbf{Framework Integration}:
    \begin{itemize}
        \item Adapt OpenR's RL suite to enhance reasoning for IT operations within AIOpsLab.
        \item Align the Q $\rightarrow$ \{R\} $\rightarrow$ A structure with IT operations workflows, building on proven successes in mathematical reasoning.
    \end{itemize}
    
    \item \textbf{Performance Optimization}:
    \begin{itemize}
        \item Evaluate the impact of OpenR's reasoning-based RL on agent performance across IT operations stages.
    \end{itemize}
    
    \item \textbf{Online Learning}:
    \begin{itemize}
        \item Implement online policy learning during inference to dynamically refine the agent's behavior and reasoning abilities as it interacts with the cloud environment.
    \end{itemize}
    
\end{enumerate}

\section{Research Methodology}

\subsection{Framework Design}
\begin{enumerate}
    \item \textbf{Agent Architecture}:
    \begin{itemize}
        \item Integrate OpenR's RL suite into AIOpsLab agents.
        \item Structure reasoning using the Q $\rightarrow$ \{R\} $\rightarrow$ A sequence for step-by-step problem-solving.
    \end{itemize}
    
    \item \textbf{Process Reward Model (PRM)}:
    \begin{itemize}
        \item Adapt OpenR's PRM to evaluate reasoning steps and outcomes for IT operations tasks.
        \item Fine-tune the PRM with IT-specific datasets to provide domain-relevant feedback.
    \end{itemize}
    
    \item \textbf{Interactive Environment}:
    \begin{itemize}
        \item Develop a simulated cloud environment for agents to practice and refine reasoning processes.
    \end{itemize}
\end{enumerate}

\subsection{Online Policy Learning}
\begin{itemize}
    \item Enable online policy learning to dynamically refine agent behavior during interactions.
    \item Utilize retrieval-augmented generation (RAG) to provide context-specific knowledge, enhancing informed decision-making.
\end{itemize}

\subsection{Dataset Preparation}
\begin{itemize}
    \item Use logs, incidents, and historical IT operations data to create a reasoning-oriented dataset.
    \item Implement \textbf{human-in-the-loop annotation} to ensure data quality and relevance.
\end{itemize}

\subsection{Evaluation Metrics}
In addition to the metrics outlined in the AIOpsLab paper, we want to propose \textbf{custom metrics} to evaluate the performance of the integrated OpenR-AIOps framework, like:
\begin{itemize}
    \item \textbf{Adaptation Speed}: Tracks the rate of performance improvement over time.
    \item \textbf{Feedback Integration}: Evaluates how effectively the agent incorporates PRM feedback to refine reasoning and actions.
\end{itemize} 

\section{Key Challenges}

\subsection{Dataset and Domain Adaptation}
Adapting OpenR's reasoning framework to IT operations tasks involves:
\begin{enumerate}
    \item \textbf{Dataset Preparation}: Unlike OpenR's math-reasoning focus, IT operations lack a well-structured dataset for process-level supervision. Creating a dataset that accurately reflects the reasoning processes involved in IT operations—such as diagnosing root causes or resolving incidents—is costly and complex.
    \item \textbf{Framework Modification}: The OpenR framework relies on Process Reward Models (PRMs) fine-tuned on task-specific datasets (e.g., OmegaPRM for math reasoning). For IT operations, the lack of pre-existing high-quality datasets for process supervision means we may need to modify the PRM design or even the reasoning structure itself (Q $\rightarrow$ \{R\} $\rightarrow$ A) to better suit IT-specific challenges.
\end{enumerate}

\subsection{Computational Costs}
Scaling RL training for reasoning-intensive tasks like IT operations is resource-intensive. OpenR's RL suite requires significant computational resources for training its models, especially when incorporating branching reasoning paths or running simulations in dynamic cloud environments.

\section{Expected Contributions}
\begin{enumerate}
    \item Integration of the OpenR framework into AIOpsLab for IT operations tasks.
    \item Development of an RL-driven methodology enhancing reasoning and decision-making in cloud environments.
    \item Establishment of a benchmark for evaluating RL-based solutions in IT operations, focusing on reasoning quality and task efficiency.
\end{enumerate}

\section{Conclusion}
Integrating OpenR into AIOpsLab offers a transformative approach to enhancing IT operations through reasoning-focused reinforcement learning. This research will develop more intelligent, adaptive, and efficient autonomous systems for managing complex cloud environments, advancing AI-driven IT management capabilities.

\begin{thebibliography}{00}
    \bibitem{b1} Y. Chen et al., “AIOpsLab: A Holistic Framework to Evaluate AI Agents for Enabling Autonomous Clouds.”
    \bibitem{b2} M. Shetty et al., “Building AI Agents for Autonomous Clouds: Challenges and Design Principles,” Jul. 31, 2024, arXiv: arXiv:2407.12165. doi: 10.48550/arXiv.2407.12165.
    \bibitem{b3} L. Zheng et al., “Judging LLM-as-a-Judge with MT-Bench and Chatbot Arena,” Dec. 24, 2023, arXiv: arXiv:2306.05685. doi: 10.48550/arXiv.2306.05685.
    \bibitem{b4} J. Wang et al., “OpenR: An Open Source Framework for Advanced Reasoning with Large Language Models,” Oct. 12, 2024, arXiv: arXiv:2410.09671. doi: 10.48550/arXiv.2410.09671.
\end{thebibliography}

\end{document}
